% =====================================================================
%  Chapter 9 — Spin, \idxw{Pole} and the \idxw{Lightcurve}
% =====================================================================

\chapter{Spin, Pole and the Lightcurve}
\label{ch:spin}

The spin state of a small body is the easiest of the missing quantities to
measure and the most consequential to get right. It is easy because a
lightcurve\idx{Lightcurve} is cheap: point a telescope at an object, watch it brighten and
dim. It is consequential because the spin state enters the probability chain
in two places and the accounting chain in none.

\section{What a lightcurve contains}

A rotating body that is not uniformly bright shows a periodic variation in
its apparent magnitude as it turns. The dominant period in that signal is the
rotation period\idx{Rotation period}, with the usual ambiguity that most small bodies present two
maxima and one minimum per rotation, so the period is usually the observed
period doubled --- unless the observer can see both maxima and minima
separately, which requires a lightcurve\idx{Lightcurve} with enough signal-to-noise and
enough sampling.

The amplitude of the variation constrains the projected extent of the body's
brightness distribution, and the shape of the variation over a rotation
constrains the body's orientation. Inverting the lightcurve\idx{Lightcurve} --- fitting a
shape model whose predicted brightness variation matches the observed one,
with the model free to rotate --- yields the rotation period\idx{Rotation period}, the direction
of the spin pole\idx{Pole} in space, and a shape model of low order.

That last product is the important one. The spin state is not a by-product
of the shape inversion; it is one of its outputs, and it is the output that
the probability chain needs.

\section{The spin state enters the probability chain twice}

The first entry is through \idxw{YORP}. The YORP\idx{YORP} effect is the torque exerted on a
body by the anisotropic re-radiation of sunlight from a rotating surface
whose normals are not uniformly distributed. The torque depends on the
surface normals, on the surface roughness distribution, and on the angle
between the spin axis and the direction of the sunlight. A shape that is
wrong changes the torque.

The second entry is through Yarkovsky. The Yarkovsky force is a recoil from
the anisotropic thermal emission of the body, and it depends on the angle
\(\gamma\) between the spin axis and the ecliptic normal through a factor
of \(\cos\gamma\). The published sensitivity table in
\tabref{tab:bottke} shows what that angle is worth: replacing the uniform
distribution of \(\cos\gamma\) with a symmetric one moves the 2036 impact
probability by a factor of 0.38, and the corresponding factor for the
spin-axis distribution is 0.23.

\section{What a shape error is worth, and where the sensitivity really is}

The 2012 study of Yarkovsky and YORP\idx{YORP} for the calibration asteroid
pseudo-Geographos, with a diameter of 2,560 metres, a bulk density of 2,500
and a semimajor axis of 1.246 astronomical units, gives the published drift
rates and the scaling laws that carry them:

\begin{equation}
\frac{d a}{d t} \propto \frac{1}{\rho_b D},
\qquad
\frac{d \omega}{d t} \propto \frac{1}{\rho_b D^{2}}.
\label{eq:rg-scaling}
\end{equation}
The size dependence is explicit and it is strong. For the same bulk density,
a halving of the diameter doubles the Yarkovsky drift rate and quadruples
the YORP\idx{YORP} torque.

The same study's sensitivity result is the one that prevents an
over-statement of the case for shape measurement, and it should be quoted
wherever the case is made. Gross shape differences move Yarkovsky by only
about ten per cent for objects of the same size. What enhances Yarkovsky
appreciably is \emph{roughness beaming}, which can raise it by tens of per
cent and in the most favourable cases by a factor of two. YORP, by
contrast, is the quantity that is genuinely shape-sensitive. And the
published propagation is: a three per cent error in the diameter produces a
three per cent error in Yarkovsky and a six per cent error in YORP.

\begin{deeper}{What this chapter is not claiming}
The book is not claiming that a resolved shape measurement buys a large
improvement in the Yarkovsky drift. On the published evidence it does not:
the drift is governed by the size, the density and the thermal parameter,
and a better shape does not change the first two. What a resolved shape buys
is the YORP term, and the spin pole\idx{Pole}, and --- as Chapter 7 established --- the
volume that the whole accounting chain starts from. The Yarkovsky
improvement comes from the density and the diameter, not from the shape
alone, and the argument for the lunar platform has to be made on the
accounting chain and on the close encounter, not on the drift.
\end{deeper}

\section{The caveat the arithmetic imposes, stated once}

The scaling laws in \eqref{eq:rg-scaling} were calibrated on a body of
2,560 metres. The class this book concerns is 50 to 140 metres, which is two
orders of magnitude smaller, and the full Yarkovsky drift carries a thermal
parameter \(\Theta = \Gamma\omega/(\varepsilon\sigma T_{\mathrm{SSP}}^{3})\)
with \(T_{\mathrm{SSP}} \propto D^{-1/2}\), so the drift does not follow from
\eqref{eq:rg-scaling} alone. **This book therefore does not extrapolate an
absolute Yarkovsky drift across that range.** The scaling laws are used only
as ratios to a calibrated reference, with the thermal-parameter dependence
left explicit, and every absolute drift rate that appears in the text is a
published value or a stated assumption.

\section{The spin state is not stationary}

The last piece of the spin picture is that it moves, and the 2029 Apophis
encounter is the demonstration. Chapter 7 set out the published result: the
gravitational perturbation of the encounter is likely to change Apophis's
effective spin rate by a factor that could be as much as two in either
direction, and to shift its pole\idx{Pole} by ten degrees or more, with the changes
very sensitive to the close-approach attitude and the mass distribution.

For the defender this has a sharp consequence. The spin state is an input to
the Yarkovsky drift, and the drift is an input to the along-track position
at a long horizon. If the spin state changes at the encounter, then a
measurement of the spin state taken \emph{before} the encounter is a
measurement of the wrong quantity, and the physical error it carries is
propagated forward along with everything else.

\begin{keytakeaways}
\item A lightcurve yields the period, the pole and a low-order shape, and it
is the cheapest of the missing measurements.
\item The spin state enters the probability chain through YORP, which is
genuinely shape-sensitive, and through the \(\cos\gamma\) factor in
Yarkovsky.
\item Gross shape differences move Yarkovsky by only about ten per cent; the
sensitivity is in YORP and in the volume.
\item The spin state is not stationary: the 2029 encounter is expected to
change Apophis's spin rate by up to a factor of two and its pole by ten
degrees or more.
\end{keytakeaways}
